\section*{Sources \& Limitations}
\addcontentsline{toc}{section}{Sources \& Limitations}

\subsection*{Primary sources consulted}
9月17日のOpenAIの法律業務向け構成、9月15日のGoogleの音声対話の新系列、9月17日のAnthropicの生命科学制度、9月のAmodei論考（9月12日づけの二次記録で日付づけ）、9月17日のAnthropic研究所発信、9月18日のAnthropicとAccentureの提携、9月15日のTypeSafeのSystem OneとJev、9月16日のCognition Code Scans、9月16日のCharacter.ai画像一家、9月16日のPixAI Tsubaki.3を範囲として読んだ。本号では全文PDFやIOC、ベンチマーク手法・図表の詳細までは扱わず、その範囲を超える主張は行わない。

\subsection*{Temporal boundary}
通常の範囲は2026-09-11 18:00から2026-09-18 18:00（いずれもAmerica/New\_York、終了時刻は含まない）。締め切り後の速報2件は通常の集計には含めない。論考の日付は月確度に二次記録の裏づけを添えて押さえ、時刻レベルの断定はしない。Jevの公開時刻は日単位で押さえ、時刻レベルの断定はしない。提携発表の時刻は頁上にないため、日単位で扱う。

\subsection*{Vendor-reported boundary}
測定値や比較、コスト、スループット、加入見立ての説明は、ベンダーやモデルカード、パートナー連携の報告として受け止め、独立した再現とは切り分ける。Valsの54.0\%や音声の82.6といった値はベンダー側・引用元側の報告であり、抽出や判定の妥当性は本号の検証範囲外である。研究所発信の計測は8月時点の切り取りであり、業界全体や9月現在の値に広げない。提携の10億ドルの額は相互の見立てであり、監査済みの支出ではない。画像の並べ比べや品質の形容はベンダー側の提示である。

\subsection*{Community observation boundary}
公開投稿の台帳（25件の直接URL：通常23件、速報2件）は関心の方向を示す背景であり、仕様やベンチマーク、日付、利用条件、価格、ライセンス、アーキテクチャ、公開事実の根拠にはならない。行単位の台帳を正とし、要約文の数えに頼らない。

\subsection*{What was not elevated}
噂の段階の話題は今週の技術的な出来事として取り上げない。二次的な言及は技術的な重みを持たせず、本号の主張には載せない。窓内に起きた運用の切り替わりは、新しいモデルの公開としては扱わず、一次資料の確認が開いたままの文脈にとどめる。
